\documentclass[conference]{IEEEtran}
\IEEEoverridecommandlockouts

\usepackage{graphicx}
\usepackage{booktabs}
\usepackage{amsmath}
\usepackage{cite}
\usepackage{url}

\title{Evidence-Driven Security Observability in an OpenStack–Kubernetes Student Cyber Range}

\author{\IEEEauthorblockN{Student A, Student B, Student C, Student D}
\IEEEauthorblockA{Centre for High Performance Computing / SCC26 Student Project\\
South Africa\\
\{student1, student2, student3, student4\}@example.invalid}
\and
\IEEEauthorblockN{Mentor Name}
\IEEEauthorblockA{Mentor / Affiliation\\
mentor@example.invalid}}

\begin{document}
\maketitle

\begin{abstract}
% 100--150 words. State: problem/context; what you built/tested; method; one or two
% evidence-backed findings; why the result matters. Do not put literature review here.
This paper presents ...
\end{abstract}

\begin{IEEEkeywords}
OpenStack, Kubernetes, GitOps, reproducibility, % add 2--4 project-specific keywords
\end{IEEEkeywords}

\section{Introduction}
% ~0.5 page.
% 1) What problem are you studying?
% 2) Why does it matter?
% 3) What gap/question does this student project address?
% 4) End with 2--4 concrete contributions, not promises.

\section{System Architecture and Context}
% ~0.5--0.7 page. Describe only the components needed to understand your experiment.
% Include one readable architecture figure. Separate Sebowa OpenStack/Kubernetes from
% externally allocated A100/H200 resources when applicable.

\begin{figure}[t]
  \centering
  % \includegraphics[width=\columnwidth]{figures/architecture.pdf}
  \fbox{\parbox[c][30mm][c]{0.92\columnwidth}{\centering Replace with architecture figure}}
  \caption{System architecture and experimental boundary.}
  \label{fig:architecture}
\end{figure}

\section{Methodology}
% ~0.7--1.0 page.
% Define reproducible procedure, versions, resources, workload/scenario, variables,
% controls, metrics and success/failure criteria. State what was held constant.

\subsection{Platform and Software Baseline}
% Git SHA(s), immutable image digest(s), Kubernetes/OpenStack resources, relevant versions.

\subsection{Experimental Procedure}
% Exact experiment sequence and number of repeats where meaningful.

\subsection{Metrics and Evidence}
% Name metrics/log sources and how timestamps/results were correlated.
% Week-5 evidence pack: one directory per run under project/runs/<run-id>/ (layout in
% project/runs/README.md). T_detect, T_identify, T_contain and T_recover are computed by
% project/tools/w5run.py from the run manifest's milestones, measured from scenario start;
% a milestone that was not reached is reported as such, never imputed.

\section{Results}
% ~1--1.3 pages. Present measured observations before interpretation.
% Use compact tables/plots. Do not report fake placeholders in the final paper.

\begin{table}[t]
\centering
\caption{Example result structure (replace with real measurements).}
\label{tab:results}
\begin{tabular}{lll}
\toprule
Configuration & Metric & Measured value \\
\midrule
A & metric-1 & TBD \\
B & metric-1 & TBD \\
\bottomrule
\end{tabular}
\end{table}

% Week-5 detection/response results. Generated by project/tools/w5paper.sh from the
% committed run directories into figures/week5/; do not edit those files by hand.
% Until a real run exists the paper says so instead of showing placeholder numbers.
\begin{table*}[t]
\centering
\caption{Week-5 detection and response times per run (from scenario start). FP/FN: false
positives/negatives; Gaps: evidence gaps. ``not reached'' means the milestone did not occur.}
\label{tab:week5-metrics}
\IfFileExists{figures/week5/metrics.tex}{%
  \resizebox{\textwidth}{!}{\input{figures/week5/metrics.tex}}%
}{%
  \fbox{\parbox{0.6\textwidth}{\centering No Week-5 run recorded yet
  (run \texttt{project/tools/w5paper.sh}).}}%
}
\end{table*}

\begin{figure*}[t]
  \centering
  \IfFileExists{figures/week5/phases.pdf}{%
    \includegraphics[width=0.75\textwidth]{figures/week5/phases.pdf}%
  }{%
    \fbox{\parbox[c][20mm][c]{0.6\textwidth}{\centering No Week-5 run recorded yet}}%
  }
  \caption{Detect, identify, contain and recover phases per Week-5 run.}
  \label{fig:week5-phases}
\end{figure*}

\section{Discussion and Limitations}
% Explain why results behaved as observed; relate to architecture/resource constraints;
% discuss measurement uncertainty, failed tests, threats to validity and what cannot be
% concluded from this small POC.

\section{Reproducibility}
% Short but explicit: repository revision, image digests, protected inputs, rebuild test,
% what was automated, what remained manual, and whether the key result was reproduced.

\section{Conclusion and Future Work}
% 1 short paragraph: answer the research question at the level supported by evidence.
% Then name the most defensible next step(s), not a wishlist.

\section*{Acknowledgment}
% Acknowledge CHPC/NICIS/CSIR, mentors, infrastructure providers and collaborators as appropriate.

% Project-specific minimum evidence ideas retained as author comments:
% - At least one approved scenario producing both host and network evidence.
% - Wazuh and Suricata detections correlated on a common timeline.
% - A human-reviewed ACP/Hermes incident explanation grounded in retained evidence.
% - Documented false positives, false negatives and telemetry gaps.
% - A clean reset/replay path that does not require attacking any system outside the authorised ranges.

\bibliographystyle{IEEEtran}
\bibliography{references}

\end{document}
